% Related work and abstract, standalone v2 paper.
% Citation keys are placeholders. Reuse whatever keys v1's .bib already defines
% and add the ones marked NEW.

\begin{abstract}
Sycophancy in language models is well documented as a behaviour but poorly localised as a
mechanism. We ask whether capitulation under user pushback, the abandonment of a correct
answer when a user objects, corresponds to a linear direction in the residual stream, and
whether such a direction emerges with scale. Across eight instruction-tuned models spanning
0.5B to 14B parameters in two families, under a pre-registered protocol with
leave-one-question-out cross-validation, within-question permutation nulls, and a positive
control, we find no usable direction below 7B and a direction that clears our validation
gate at 7B, 8B and 14B. Ablating that direction in Qwen2.5-7B reduces the rate at which the model reverses a correct answer by 4.8 percentage points, roughly a quarter in relative
terms, while the identical intervention on a gate-failing direction and on a marginally
passing one produces no effect. Probe quality, not statistical significance, predicts causal
relevance. Behaviourally, pushback destroys three to five times more correct answers than it
repairs at 1B and reaches parity by 7B, falsifying our own pre-registered prediction, and a
sign-reversing crossover between pushback types survives an order of magnitude of scale and
marginalisation over five paraphrases per type. We additionally quantify four measurement
hazards, including a numerical-precision effect that changes roughly one episode label in
ten.
\end{abstract}

